\section{Absolute convergence of the doubled Eisenstein series at zero}

% BEGIN C75E1
\paragraph{C75E1. Representatives with a common denominator.}
Write $E$ for the CM field, $\mathcal O$ for its ring of integers,
$\Psi$ for its CM type and $s=|\Psi|$, so $[E:\mathbb Q]=2s$.
The split matrix $J$ has diagonal blocks zero and off-diagonal blocks
$\epsilon I,-\epsilon I$, with $\epsilon\in\{1,-1\}$. Each index
$U$ in O74E1 is a totally isotropic $r$-dimensional row space in
$E^{2r}$. There is a single positive integer $c$ and for every such
$U$ a matrix $g_U$ preserving $J$, with bottom space $U$ and
$c g_U$ integral.

Indeed $P=U\cap\mathcal O^{2r}$ has rank $r$ and is saturated.
The quotient $\mathcal O^{2r}/P$ is torsion-free, hence projective
over the Dedekind ring $\mathcal O$. In particular $P$ is a direct
summand. By the Steinitz classification of projective modules,
$P\simeq\mathcal O^{r-1}\oplus\mathfrak a$ for one of a fixed
finite list of integral ideal class representatives. For each such
ideal choose a nonzero $b\in\mathfrak a$. The finite quotient
$\mathfrak a/b\mathcal O$ is killed by a positive integer. A common
multiple $c_0$ for the list gives a free submodule $F\subset P$ of
rank $r$ containing $c_0P$, by using $b\mathcal O$ in the last slot.
Let $W$ have as rows a basis of $F$. The coordinate forms dual to
these rows take $P$ into $c_0^{-1}\mathcal O$. Extend them by the
integral projection onto $P$. Conjugating the coordinate forms and
using $J^{-1}$ gives a matrix $V$ with $c_0V$ integral and
$VJW^*=\epsilon I$. Explicitly, if the extended duals form a matrix
$D$ of columns with $WD=I$, take $V=\epsilon D^*J^{-1}$.

The matrix $K=VJV^*$ is skew-Hermitian. Replacing $V$ by
$V'=V+(\epsilon/2)KW$ preserves the mixed pairing and kills the
top Gram block: $WJV^*=-\epsilon I$, $WJW^*=0$, and the two cross
terms each equal $-K/2$. The stacked rows $(V';W)$ thus have Gram
matrix $J$. Their matrix is invertible, so it also preserves $J$ in
the source's column convention. All entries belong to
$(2c_0^2)^{-1}\mathcal O$. This gives $c=2c_0^2$, independent of
$U$. For $r=0$ there is one index and one empty representative.
% END C75E1

% BEGIN C75E2
\paragraph{C75E2. Counting by short integral generators.}
Let $W=(C\ D)$ be the bottom rows of one of these representatives,
and put $H(U)=\prod_{\tau\in\Psi}|\det(iC_\tau+D_\tau)|$. For
$n>4r$ we have $\sum_U H(U)^{-n}<\infty$. We give the count since
it is here that the strict inequality on $n$ is needed.

Embed the rows of $cW$ at all CM places into
$\prod_{\tau\in\Psi}\mathbb C^{2r}$. Their $\mathcal O$-span $L$
is a real lattice of rank $2sr$ in its span. Writing $G_\tau=
W_\tau W_\tau^*$, its Euclidean covolume in that span is
\[
 \operatorname{covol}(L)=C_{E,r}\,c^{2sr}
      \prod_{\tau\in\Psi}\det G_\tau
   = C_{E,r}\,c^{2sr}H(U)^2.
\]
For the first equality take the Minkowski lattice $\mathcal O^r$ as
domain; the map on rows at $\tau$ has real volume multiplier
$\det(W_\tau W_\tau^*)$, the product of its squared complex singular
values. For the second, $C_\tau D_\tau^*=D_\tau C_\tau^*$ by isotropy,
so $(iC_\tau+D_\tau)(iC_\tau+D_\tau)^*=G_\tau$.

Write $d=2s$ and $\lambda_j$ for the Euclidean successive minima of
$L$. Choose $r$ vectors $w_1,\ldots,w_r$ successively over $E$ with
$\|w_{i+1}\|\le\lambda_{di+1}$ for $0\le i<r$. Indeed among
$di+1$ independent rational vectors of this length one lies outside
the $E$-span of the first $i$, a rational space of dimension $di$.
Monotonicity of the minima and Minkowski's second theorem give
\[
 \prod_{i=1}^r\|w_i\|^{2s}
 \le\prod_{j=1}^{2sr}\lambda_j
 \le C_{2sr}\operatorname{covol}(L) .
\]
These are nonzero integral vectors in $\mathcal O^{2r}$ and their
$E$-span recovers $U$. Choose one ordered tuple per $U$. The resulting
map is injective; the preceding bounds give
$H(U)^{-n}\le C\prod_i\|w_i\|^{-sn}$ with a constant independent
of $U$. The lattice $\mathcal O^{2r}$ has real rank $4sr$. In the
shell $2^j\le\|w\|<2^{j+1}$ it has $O(2^{4srj})$ points by lattice
packing. The finitely many nonzero points of length below one cause
no difficulty. Thus $\sum_{0\ne w\in\mathcal O^{2r}}\|w\|^{-sn}$
converges when $sn>4sr$. Raising this convergent sum to the $r$th
power bounds the entire sum over $U$.
% END C75E2

% BEGIN C75E3
\paragraph{C75E3. A uniform bound for the finite test value.}
Fix a positive level $R$, an even $n$ and a native Hermitian
Schwartz--Bruhat test $f$ of level $R$ on $\operatorname{Mat}_{n,r}(E)$.
It is periodic modulo $R\mathcal A$ and supported on
$R^{-1}\mathcal A$, where $\mathcal A=\operatorname{Mat}_{n,r}(\mathcal O)$.
For all split unitary $g$ with $cg$ integral there is a bound
\[
         |(\omega(g) f)(0)|\le
         (2cR^2)^{[E:\mathbb Q]nr/2}\,\|f\|_\infty .
\]
We use the native Weil action's preservation of the finite inner
product and its Heisenberg covariance. If $d,b\in2R\mathcal A$,
the displacement $H(d,b)$ fixes $f$. Translation by $d$ is a
period. On the support the trace in its phase is
$\operatorname{Tr}_{E/\mathbb Q}\sum b(2\bar x-\bar d)$, an even
integer, so the exponential with coefficient $\pi i$ is one.
Outside the support the translation preserves vanishing.
For $t\in2cR\mathcal A$, the two blocks of $(t,0)g$ lie in
$2R\mathcal A$. Weil covariance transports this fixed displacement
to $H(t,0)$ on $\omega(g)f$. The transformed test is therefore
constant on $2cR\mathcal A$. Integrating its squared modulus there
and using preservation of the inner product gives
\[
 |(\omega(g)f)(0)|^2\operatorname{vol}(2cR\mathcal A)
 \le\|f\|_2^2\le\operatorname{vol}(R^{-1}\mathcal A)\|f\|_\infty^2.
\]
For the source's rational Haar functional the ratio of the two
volumes is the index $(2cR^2)^{[E:\mathbb Q]nr}$, as follows by
partition into cosets. This proves the bound with the original
normalization.
% END C75E3

% BEGIN C75E4
\paragraph{C75E4. The stable-range majorant at zero elliptic arguments.}
For $R>0$, $\delta>0$, even $n>4r$, there is a constant
$C(E,\Psi,n,r,R,\delta)$ such that every level-$R$ test $f$ and every
tuple of periods with $\operatorname{Im}Z_\tau\ge\delta I$ satisfy
\[
 \sum_U |\mathrm{eisTerm}(g_U,f;Z,0,0)|
             \le C(E,\Psi,n,r,R,\delta)\,\|f\|_\infty .
\]
In particular the absolute convergence input in O74E4 and O74E5 is
available at every allowed period for the actual doubled finite tests.

At one embedding take the bottom blocks $(C,D)$ of any representative.
Write $B(Z)=(CZ+D)^{-1}C$. The denominator is invertible on the
Hermitian half space. For example if $(CZ+D)^*v=0$, then with
$u=C^*v$ one has $D^*v=-Z^*u$. Isotropy $CD^*=DC^*$ forces
$u^*(Z-Z^*)u=0$, whence $u=0$ and $D^*v=0$. Full bottom-row rank
then gives $v=0$. By O74E4, $-\operatorname{Im} B(Z)=
B(Z)(\operatorname{Im} Z)B(Z)^*\ge0$ and $\|B(Z)\|\le d^{-1}$
when $\operatorname{Im} Z\ge dI$. Along $Z(t)=X+iY(t)$ with
$X=X^*$ and $Y'(t)\ge0$ the logarithmic modulus derivative is
\[
 \frac{d}{dt}\log|\det(CZ(t)+D)|
    =\operatorname{Re}\operatorname{tr}(B(Z(t))\,iY'(t))
    =\operatorname{tr}((-\operatorname{Im}B(Z(t)))Y'(t))\ge0 .
\]
Therefore at a fixed $X$ the modulus for $Y\ge\delta I$ is bounded
below by its value for $Y=\delta I$. When $\|X\|\le K$, join $iI$
to $X+i\delta I$ by the line segment. The imaginary part stays at
least $\min(1,\delta)I$. Integrating the bound
$|\operatorname{tr}(B Z')|\le r\|B\|\|Z'\|$ along the segment gives
\[
 |\det(C(X+iY)+D)|\ge \exp\!\left(
   -r(K+|1-\delta|)/\min(1,\delta)\right)|\det(iC+D)|.
\]

To bound all real parts at once write $X_\tau=(Z_\tau+Z_\tau^*)/2$.
There is an integral Hermitian matrix $\beta$ for which every
$\|X_\tau-\tau\beta\|\le K$, with $K$ depending only on $E,r$.
Indeed the off-diagonal entries can be rounded in a fundamental
parallelepiped of the full Minkowski lattice $\mathcal O$ in
$\mathbb C^s$ and the diagonal entries in one for the real integers
$\mathcal O_{E^+}$ in $\mathbb R^s$. Fill the other triangle by
conjugation. Put $t_\beta=\left(\begin{smallmatrix}I&\beta\\0&I\end{smallmatrix}\right)$.
Right multiplication $U\mapsto U'=Ut_\beta$ permutes the indices.
Use $g_{U'}t_{-\beta}$ as a representative of $U$. It still has
denominator dividing $c$. If $(C,D)$ are the bottom blocks of
$g_{U'}$, the denominator at $Z$ is $C(Z-\beta)+D$. The preceding
comparison, at all $s$ places, bounds its inverse determinant factor
to power $n$ by a fixed constant times $H(U')^{-n}$. At zero elliptic
arguments the numerator is $(\omega(g_{U'}t_{-\beta})f)(0)$ and the
Jacobi exponential is one. By representative independence in O74E3
and the uniform bound of C75E3 the absolute term is at most
$C'\|f\|_\infty H(U')^{-n}$. C75E2 gives the stated estimate.

For the companion theta bound at zero, sum the ordinary Gaussian
$\exp(-2\pi\delta\|u\|^2)$ over the fixed denominator lattice
$R^{-1}\mathcal A$ at the CM places. This also bounds the absolute
theta series by a constant times $\|f\|_\infty$. Each doubled test
in O74E3 has some positive level by its defining Schwartz--Bruhat
property. Choosing that level and a common positive lower eigenvalue
for its period tuple gives precisely the convergence used there.
% END C75E4
